\section{Conclusiones}
\label{sec:conclusiones}

El trabajo desarrolló y verificó un flujo completo de apoyo a la decisión para el dimensionamiento del sector de preparación de pedidos online de un supermercado, integrando el pronóstico horario por media estacional de ocho semanas, la verificación analítica mediante el modelo $M/M/c$ y la simulación de eventos discretos con números aleatorios comunes en SimPy.
Todos los objetivos planteados se alcanzaron de manera reproducible a través de una tubería automatizada cuyos parámetros, semillas y resultados quedaron registrados y contrastados sobre un período de prueba reservado.

\subsection{Recomendaciones}
\label{sec:recomendaciones}

Para cumplir el objetivo operativo de preparar al menos el 90~\% de los pedidos online en un plazo máximo de 30~minutos desde su ingreso, se recomienda asignar la dotación de preparadores según el escenario diario de demanda pronosticado para el sistema modelado:
\begin{itemize}
  \item \textbf{Jornadas de demanda baja ($\approx 159$~pedidos/día, como los domingos):} asignar $c^* = 3$~preparadores, con lo cual se alcanza un cumplimiento estimado del 92,0~\% (90,6~\% en los días de prueba), una espera media en fila de 5,52~minutos y una utilización del 65,5~\%.
  En fechas en que se prevea un crecimiento superior al 10~\% sobre el pronóstico base o ante campañas promocionales, se sugiere ampliar la dotación a $c = 4$~preparadores debido al reducido margen de holgura de este escenario.
  \item \textbf{Jornadas de demanda media ($\approx 185$~pedidos/día, como los jueves):} asignar $c^* = 4$~preparadores, logrando un 97,4~\% de pedidos a tiempo (100,0~\% en los días de prueba), una espera media de 2,87~minutos y una utilización del 58,4~\%, con holgura suficiente para absorber subestimaciones del pronóstico de hasta un 10~\%.
  \item \textbf{Jornadas de demanda alta ($\approx 300$~pedidos/día, como los sábados):} asignar $c^* = 6$~preparadores, alcanzando un cumplimiento del 97,5~\% (98,1~\% en los días de prueba), una espera media de 2,82~minutos y una utilización del 63,2~\%.
  \item \textbf{Contingencia ante mayores tiempos de preparación:} si una medición del tiempo de armado determinara que su media supera los 10~minutos supuestos y se aproxima a 12~minutos por pedido ($\times 1{,}20$), debe adicionarse un preparador en cada escenario ($c^* = 4$, 5 y 7).
\end{itemize}

\subsection{Limitaciones}
\label{sec:limitaciones}

Las conclusiones del estudio deben interpretarse dentro de los límites de los datos y supuestos adoptados:
\begin{itemize}
  \item \textbf{Naturaleza sintética de los datos:} debido al bajo volumen y a la ausencia de patrón horario del histórico del ERP, las llegadas se generaron mediante un proceso de Poisson no homogéneo calibrado para un supermercado regional (unos 200~pedidos por día), por lo que los valores numéricos reflejan el comportamiento del sistema modelado y no mediciones de campo de un comercio particular.
  \item \textbf{Supuesto sobre la duración del servicio:} al no existir en el ERP el instante de finalización del armado, el tiempo de preparación se modeló como una función lineal determinística de la cantidad de líneas ($s(\ell) = b + m\ell$), sin variabilidad adicional para pedidos de igual tamaño ni diferencias de productividad entre empleados.
  \item \textbf{Dotación fija por jornada:} el modelo evalúa plantillas constantes durante las 13~horas de atención y no contempla el escalonamiento de turnos de entrada y salida dentro del mismo día.
  \item \textbf{Alcance acotado al armado:} se asumió disponibilidad plena de inventario y no se modelaron quiebres de stock, sustituciones de productos, recorrido físico en góndola ni la etapa posterior de despacho a domicilio.
\end{itemize}

\subsection{Extensiones}
\label{sec:extensiones}

A partir de la infraestructura computacional desarrollada, el estudio admite cuatro extensiones directas:
\begin{itemize}
  \item \textbf{Registro de tiempos de preparación en el ERP:} incorporar en la tabla de pedidos web la marca temporal de fin de preparación para estimar empíricamente la distribución del tiempo de servicio según el número de líneas del pedido.
  \item \textbf{Turnos escalonados:} reemplazar la dotación constante diaria por una capacidad variable por franja horaria $c_h$, de modo de cubrir los picos de las 11 y las 19~horas con menor tiempo ocioso en las horas de baja demanda.
  \item \textbf{Prioridades y preparación por lotes:} evaluar disciplinas de cola con prioridad según la modalidad de entrega (retiro en sucursal o envío) y estrategias de preparación simultánea de varios pedidos.
  \item \textbf{Integración con el tablero gerencial del ERP:} conectar el paquete \texttt{pedidos} al tablero del sistema desarrollado en Sistemas~III para actualizar semanalmente el pronóstico y sugerir la dotación de cada día de la semana entrante.
\end{itemize}
